Clinical early warning systems are evaluated almost exclusively by discrimination, yet their failure in deployment is operational rather than statistical: intensive care units generate alarm volumes far beyond what clinical staff can absorb, and the resulting alarm fatigue erodes the very safety the systems are meant to provide. This paper treats the alarm budget, defined as the tolerated number of alarms per patient-day, as a first-class design constraint rather than a post hoc threshold adjustment. We formalise early warning as a stateful alarm policy whose refractory window and escalation tier are explicit control variables, and we show that the resulting temporal dependence invalidates timestep-level conformal calibration. We therefore develop a patient-level blocked conformal risk control procedure that restores finite-sample validity for dependent alarm sequences and yields an upper bound on the probability that the deployed population alarm rate exceeds the operational budget. The method is instantiated on irregularly sampled intensive care trajectories using two open, multi-hospital cohorts, and is evaluated under cross-hospital transfer, prevalence shift, and observation sparsification. The AI contribution is a calibration procedure that provides distribution-free control of alarm rate for policies with internal state; the engineering contribution is an operating-point selection mechanism that lets a monitoring system be specified by its alarm budget and then certified against it. The evaluation protocol compares all methods at matched alarm budgets rather than at matched thresholds, reports detection lead time and net benefit alongside discrimination, and quantifies the empirical budget violation rate over repeated calibration splits.
